El objetivo de este artículo es analizar la integración de patrones de diseño, el modelo-vista-controlador (MVC) y la programación orientada a objetos (POO) en el desarrollo de proyectos de software, destacando su potencial para mejorar la organización y mantenibilidad del sistema. Para alcanzar este propósito, se adoptó un enfoque metodológico aplicado en un proyecto académico, en el cual se implementaron dichas prácticas como base en la construcción del sistema. Aunque no se implementó de manera estricta la metodología Scrum, se incorporaron algunas de sus prácticas, como las reuniones diarias, que permitieron evaluar sus beneficios en la gestión ágil y la comunicación entre los integrantes del equipo. Los resultados obtenidos evidencian que la combinación de MVC y POO facilita la separación de responsabilidades y la reutilización del código, mientras que la integración parcial de Scrum potencia la organización colaborativa. Se concluye que este tipo de integración representa un ejemplo de prácticas modernas en el desarrollo de software, aplicables tanto en contextos académicos como profesionales.
